\section{Evaluation methodology and simulated results}\label{sec:eval}

In a simulation with stated assumptions, micro-batching makes server-side execution about
eight times cheaper at a cost of a few milliseconds, and Markov prefetch hides most network
latency but multiplies compute. These numbers come from a model, not a deployment;
\cref{sec:realeval} lists what a real evaluation must measure.

\subsection{Setup}
\code{simulate.py} (\cref{app:simulate}) uses SimPy \cite{simpy} for a discrete-event model
of one protected feature. Requests arrive as a Poisson process at 2,000 per second. A
batched kernel call on \(b\) items costs \(4.0 + 0.15\,b\) ms, with at most 64 items per
batch; the batching timer starts when the first request enters an empty queue. The GPU pool
is sized from the expected batch size so that utilisation stays at or below 70\%. Each
configuration simulates 20\,s (about 40,000 requests), discarding the first second. For
the prefetch experiment, sessions of 20 steps are walks on a skewed first-order Markov chain
over 12 features; a model fitted on the first 1,000 of 3,000 sessions is evaluated on the
remaining 2,000. A miss costs the network round trip (40\,ms) plus one unbatched kernel call
(4.15\,ms); a prefetched local hit costs 5\,ms. Speculative runs that are never used count
as extra compute.

\subsection{Results}
\pgfplotstableread[col sep=comma]{code/results_batching.csv}\batching
\pgfplotstableread[col sep=comma]{code/results_prefetch.csv}\prefetch

\paragraph{Batching window.} \Cref{fig:batching} plots GPU time per 1,000 requests against
server-side 95th-percentile (p95) latency for windows from 0 to 32\,ms, and
\cref{tab:batching} gives the full results. Without batching, each request pays the full
fixed cost, so 1,000 requests take 4,150 GPU-ms on a pool of 12 GPUs. A 5\,ms window
gathers 11 requests per batch on average and cuts this to 513 GPU-ms on 2 GPUs, while
server p95 rises from 5.0 to 10.9\,ms. Beyond 20\,ms the fixed cost is already spread over
about 40 requests, so cost falls only from 247 to 215 GPU-ms while p95 rises from 29.8 to
42.7\,ms; at 32\,ms most batches reach the 64-item cap.

\begin{figure}[htbp]
\centering
\begin{tikzpicture}
\begin{axis}[
    width=0.86\linewidth, height=6.2cm,
    font=\sffamily\footnotesize,
    xlabel={Server p95 latency (ms), excluding network round trip},
    ylabel={GPU-ms per 1,000 requests},
    xmin=0, xmax=46, ymin=0, ymax=4600,
    ytick={0,1000,2000,3000,4000},
    scaled y ticks=false,
    yticklabel style={/pgf/number format/.cd, 1000 sep={,}},
    grid=major, grid style={draw=black!8},
    axis line style={draw=black!40}, tick style={draw=black!40},
    nodes near coords={\(W{=}\pgfplotspointmeta\)\,ms},
    nodes near coords style={font=\sffamily\scriptsize, text=textgray, anchor=south west, inner sep=2pt},
    point meta=explicit symbolic]
  \addplot[color=boxgray, mark=*, mark size=1.6pt, line width=0.8pt]
      table[x=p95_ms, y=gpu_ms_per_1k, meta=window_ms] {\batching};
  \addplot[only marks, mark=*, mark size=3pt, color=accent, nodes near coords={}]
      table[x=p95_ms, y=gpu_ms_per_1k, meta=window_ms,
            restrict expr to domain={\thisrow{window_ms}}{5:5}] {\batching};
\end{axis}
\end{tikzpicture}
\caption{A 5\,ms batching window (blue) cuts GPU time per request about eightfold for
6\,ms more p95 latency; gains flatten past 20\,ms. \code{simulate.py}, SimPy model,
2,000 req/s, assumed kernel costs (not measured). \(W = 0\) is no batching.}
\label{fig:batching}
\end{figure}

\begin{table}[htbp]
\centering
\caption{Batching-window results from \code{simulate.py} (server-side latency only).}
\label{tab:batching}
\small
\pgfplotstabletypeset[
  columns={window_ms,gpus,mean_batch,gpu_ms_per_1k,utilisation,p50_ms,p95_ms,p99_ms},
  columns/window_ms/.style={column name={\(W\) (ms)}, column type={r}},
  columns/gpus/.style={column name={GPUs}, column type={r}},
  columns/mean_batch/.style={column name={Mean batch}, fixed, fixed zerofill, precision=1, column type={r}},
  columns/gpu_ms_per_1k/.style={column name={GPU-ms / 1k req}, fixed, fixed zerofill, precision=1,
                                 1000 sep={{,}}, column type={r}},
  columns/utilisation/.style={column name={Util.}, fixed, fixed zerofill, precision=2, column type={r}},
  columns/p50_ms/.style={column name={p50 (ms)}, fixed, fixed zerofill, precision=1, column type={r}},
  columns/p95_ms/.style={column name={p95 (ms)}, fixed, fixed zerofill, precision=1, column type={r}},
  columns/p99_ms/.style={column name={p99 (ms)}, fixed, fixed zerofill, precision=1, column type={r}},
  every head row/.style={before row=\toprule, after row=\midrule},
  every last row/.style={after row=\bottomrule}]{\batching}
\end{table}

The knee of this curve depends on traffic: a feature with 20 requests per second will
rarely fill a batch within 5\,ms, so low-traffic features gain far less. This supports the
lifecycle policy in \cref{sec:lifecycle}: features that cannot batch well carry higher
serving pressure \(S\) and become candidates for promotion sooner.

\paragraph{Prefetch depth.} \Cref{fig:prefetch} varies the number of speculated next
features \(k\). With \(k = 2\), 68\% of next steps are already in the client cache and mean
perceived latency falls from 44.1 to 17.4\,ms, but speculative work adds 132\% to demand
compute. The batching savings in \cref{fig:batching} are large enough to pay for
\(k = 1\) or \(2\); \(k = 4\) adds 309\% for a further 8.7\,ms saving.

\begin{figure}[htbp]
\centering
\begin{tikzpicture}
\begin{axis}[
    width=0.86\linewidth, height=6.2cm,
    font=\sffamily\footnotesize,
    xlabel={Extra speculative compute (\% of demand compute)},
    ylabel={Mean perceived latency (ms)},
    xmin=-10, xmax=345, ymin=0, ymax=52,
    xticklabel={\pgfmathprintnumber{\tick}\%},
    grid=major, grid style={draw=black!8},
    axis line style={draw=black!40}, tick style={draw=black!40},
    visualization depends on={100*\thisrow{hit_rate} \as \hits},
    nodes near coords={\(k{=}\pgfplotspointmeta\), \pgfmathprintnumber[fixed,precision=0]{\hits}\% hits},
    nodes near coords style={font=\sffamily\scriptsize, text=textgray, anchor=south west, inner sep=3pt},
    point meta=explicit symbolic]
  \addplot[color=boxgray, mark=*, mark size=1.6pt, line width=0.8pt]
      table[x=extra_compute_pct, y=mean_latency_ms, meta=k] {\prefetch};
  \addplot[only marks, mark=*, mark size=3pt, color=accent, nodes near coords={}]
      table[x=extra_compute_pct, y=mean_latency_ms, meta=k,
            restrict expr to domain={\thisrow{k}}{2:2}] {\prefetch};
  \draw[dashed, draw=textgray] (axis cs:-10,40) -- (axis cs:345,40)
      node[pos=0.98, above left, font=\sffamily\scriptsize, text=textgray] {Network round trip, 40\,ms};
\end{axis}
\end{tikzpicture}
\caption{Prefetching the two most likely next steps (blue) cuts mean wait from 44.1 to
17.4\,ms for 132\% extra speculative compute. \code{simulate.py}, 3,000 synthetic sessions
over 12 features, first third used for training.}
\label{fig:prefetch}
\end{figure}

\subsection{What a real evaluation must measure}\label{sec:realeval}
The simulation fixes kernel costs, traffic and user behaviour by assumption. A deployment
study should report:
(1)~per-feature p50, p95 and p99 end-to-end latency across real networks, since tail latency
dominates experience \cite{dean2013tail};
(2)~GPU cost per 1,000 invocations with real traffic mixes and multi-tenant batching;
(3)~prefetch hit rate and wasted compute from real session logs;
(4)~a red-team exercise in which an AI-assisted reverse-engineering agent targets the client
shell, reporting what it recovers; and
(5)~an extraction study measuring how many queries a substitute model needs to reach a given
fidelity, with and without the controls of \cref{sec:security}
\cite{tramer2016stealing,juuti2019prada,orekondy2019knockoff}.
